\documentclass[10pt,a4paper]{article}
\usepackage[hmargin=1.5cm,top=1.1cm,bottom=1.1cm]{geometry}
\usepackage[T1]{fontenc}
\usepackage[utf8]{inputenc}
\usepackage{lmodern}
\usepackage{microtype}
\usepackage{booktabs}
\usepackage{graphicx}
\usepackage{xcolor}
\usepackage{tcolorbox}
\usepackage{enumitem}
\usepackage[hidelinks]{hyperref}
\usepackage{soul}
\sethlcolor{yellow!55}
\soulregister\texttt7
\soulregister\emph7
\soulregister\textbf7

\setlist{nosep,leftmargin=*}
\setlength{\parskip}{2.5pt}
\setlength{\parindent}{0pt}
\pagestyle{empty}
\newcommand{\figdir}{../output/figures}
\newcommand{\q}[1]{\textbf{\color{blue!50!black}#1}}

\begin{document}
\small

{\large\textbf{Predicting credit card churn with K-Nearest Neighbours}}\\
Ignacio Queipo de Llano \textperiodcentered{} Predictive Machine Learning in Finance, IESE MiF \textperiodcentered{} 18 November 2026\\[-6pt]
\rule{\linewidth}{0.4pt}

\begin{tcolorbox}[colback=white,colframe=black,boxrule=0.4pt,sharp corners,left=5pt,right=5pt,top=3pt,bottom=3pt]
\textbf{Bottom line.} 16.1\% of the 10,127 cardholders closed their account, so ``nobody leaves'' is 84\% accurate.
KNN only works once distances mean something: unscaled it scores 89.1\% \emph{by accident}, because three dollar columns
make up 99.6\% of the distance and happen to be predictive. The right pipeline (scale the numbers, one-hot the
categories, $k=7$ by cross-validation) reaches 90.5\% accuracy, 83\% precision and AUC 0.91. For a retention campaign,
recall matters more: with the alarm at a 20\% churn probability the model catches 81\% of churners, two in three alarms
being right.
\end{tcolorbox}

\q{Q1. Data.} \hl{10,127 customers, 23 columns; 16.07\% churn (1,627).} Dropped the two
\texttt{Naive\_Bayes\_Classifier} columns (predictions built with the target: leakage) and the ID. \hl{14 numeric, 5
categorical} (\texttt{Gender}, \texttt{Education\_Level}, \texttt{Marital\_Status}, \texttt{Income\_Category},
\texttt{Card\_Category}); 37 columns after \texttt{get\_dummies}. Stratified 80/20 split: 8,101 / 2,026.

\q{Q2. KNN, $k=5$, raw data:} \hl{89.1\% accuracy, recall 55\%.} It beats the 84\% baseline, but for the wrong reason (Q3).

\q{Q3. Scale.} \hl{The \texttt{Credit\_Limit} range (33,078) is 33,111 times that of \texttt{Avg\_Utilization\_Ratio}
(0.999).} Euclidean distance adds squared gaps, so \hl{\texttt{Credit\_Limit} and \texttt{Avg\_Open\_To\_Buy} take 93\%
of it}: KNN is really a ``similar credit line'' rule and ignores ratios, counts and dummies.

\q{Q4. Scaling.} StandardScaler: $(x-\mu)/\sigma$; MinMaxScaler: $(x-\min)/(\max-\min)$, which an outlier squashes.
\hl{Scaling all 37 columns makes things worse: 85.8\% (recall 22\%).} The dummies are the problem: a Platinum card
(rare) becomes 21.8 s.d. when it is 1, so one rare category dominates the distance. \hl{For KNN: StandardScaler on the
numeric columns only, dummies left as 0/1.}

\vspace{2pt}
\begin{minipage}[t]{0.46\linewidth}\vspace{0pt}
\includegraphics[width=\linewidth]{\figdir/q5_accuracy_vs_k.png}
\end{minipage}\hfill
\begin{minipage}[t]{0.51\linewidth}\vspace{0pt}
\q{Q5. Choosing $k$} (5-fold CV, pipeline). \hl{$k=3$ overfits: 94.8\% train vs 90.7\% validation.} The gap closes
as $k$ grows; \hl{from $k=21$ both lines fall together: underfitting} (51 neighbours average away the signal).
\hl{Best $k=7$: 91.1\% validation.} The curve is flat between 5 and 15, so the exact $k$ matters little.

\q{Q7. Curse of dimensionality.} In high dimensions all points are almost equally far apart, so the ``nearest''
neighbour is barely nearer than the rest. Uniform simulation: \hl{nearest/farthest distance ratio 0.01 in 2 dimensions,
0.52 in 37.} Remedies: fewer, better features (dropping the 23 dummies gives 91.2\%, the best model), PCA, feature
weighting, or a model that selects features (trees).
\end{minipage}

\vspace{3pt}
\q{Q6. Best model and threshold} (test set, 2,026 customers, 325 churners):

\vspace{2pt}
\centerline{\begin{tabular}{lrrrrrrr}
\toprule
 & TP & FN & FP & TN & Accuracy & Recall & Precision\\
\midrule
Always ``stays'' & 0 & 325 & 0 & 1,701 & 84.0\% & 0\% & --\\
KNN $k=5$, unscaled & 179 & 146 & 75 & 1,626 & 89.1\% & 55.1\% & 70.5\%\\
KNN $k=5$, StandardScaler on all columns & 71 & 254 & 33 & 1,668 & 85.8\% & 21.8\% & 68.3\%\\
Pipeline $k=7$, cut-off 50\% & 167 & 158 & 35 & 1,666 & 90.5\% & 51.4\% & \textbf{82.7\%}\\
Pipeline $k=7$, cut-off 20\% & 263 & 62 & 143 & 1,558 & 89.9\% & \hl{\textbf{80.9\%}} & 64.8\%\\
Numeric columns only, $k=11$ & 185 & 140 & 39 & 1,662 & \textbf{91.2\%} & 56.9\% & 82.6\%\\
\bottomrule
\end{tabular}}

\vspace{2pt}
\emph{Business reading.} \hl{Recall matters more:} a missed churner is a lost customer worth five to seven
acquisitions; a false alarm is a retention offer to someone who was staying. With $k=7$ the share of churning neighbours
can only take values in sevenths, so the 20\% cut-off means ``2 of 7 neighbours left''.

\q{Q8. Pipeline.} \hl{\texttt{ColumnTransformer} (scale numeric, one-hot categorical) + KNN, fitted on the raw training
data, gives exactly the same test predictions as the manual steps.} It is better because the scaler is learnt inside each
CV fold (no test information leaks into $\mu$, $\sigma$), the same object transforms new customers in production, and
the whole thing can be tuned and saved as one model.
\end{document}
